\section{Runge 现象与节点选择}
\label{sec:08-Numerical-Analysis/03-Interpolation-and-Approximation/02}

你有一个看起来很温顺的函数

\[
f(x)=\frac{1}{1+25x^2},\qquad x\in[-1,1].
\]

没有尖点，没有跳跃，无限阶可导，图像光滑得像个倒扣的钟。你在 \([-1,1]\) 上等距地取 \(n+1\) 个点，做一个 \(n\) 次插值多项式，让它精确穿过这些点。凭直觉，\(n\) 越大，采样越密，逼近应该越来越好——插值多项式在所有数据点上误差本来就是零，点多了，两个采样点之间的空隙也更小。

\(n=5\) 的时候，插值曲线紧贴原函数，看起来很好。\(n=9\)，中间区域仍然不错，但端点附近开始微微上扬。\(n=13\)，端点处的摆动肉眼可见。\(n=21\)，端点附近的振荡幅度已经远大于函数值本身——而这些位置离最近的采样点不过一个子区间的距离。插值曲线在所有 \(n+1\) 个数据点上精确命中，误差为零；在它们之间的区间上，却完全不可信任。

这就是 Runge 现象。它精准地粉碎了一个自然的信念：更多数据点加上更高次的多项式模型，必然给出更准确的函数逼近。数据点越多，端点越狂暴——这不是实现上的小瑕疵，而是等距节点高次插值的结构性灾难。

\subsection{问题不在多项式空间，而在插值这个动作本身}

要理解发生了什么，先回到插值的定义。记 \(I_nf\) 为在 \(n+1\) 个节点上插值 \(f\) 得到的 \(n\) 次多项式。如果数据有微小扰动——比如每个 \(y_j\) 浮动 \(\delta y_j\)——输出的多项式就会变成

\[
\delta p(x)=\sum_j\delta y_j\,\ell_j(x).
\]

最坏情况的放大倍数是

\[
\Lambda_n=\max_{x\in[-1,1]}
\sum_{j=0}^n|\ell_j(x)|,
\]

称为 Lebesgue 常数。它衡量插值算子对数据扰动的敏感度——条件数。等距节点下的 \(\Lambda_n\) 随 \(n\) 近似以指数速度增长：\(n=10\) 时约 30；\(n=20\) 时约 10,000；\(n=30\) 时超过 \(10^7\)。这意味着在 \(n=30\) 的等距插值中，数据上一个 \(10^{-16}\) 量级的扰动（双精度的舍入量级）可以在端点被放大到 \(10^{-9}\)——七个数量级的精度蒸发。

但这还只是扰动放大的故事。要理解 Runge 现象本身——即真实函数值和插值多项式之间的系统性偏离——我们需要 Lebesgue 不等式。

记 \(P_n\) 为次数不超过 \(n\) 的多项式全体。任取 \(q\in P_n\)，因为 \(q\) 在自己节点上的插值就是自己（插值算子的不动点），有 \(I_nq=q\)。于是

\[
f-I_nf=(f-q)-I_n(f-q).
\]

取一致范数 \(\|\cdot\|_\infty\)，利用插值算子对扰动的放大估计 \(\|I_ng\|_\infty\le\Lambda_n\max_j|g(x_j)|\le\Lambda_n\|g\|_\infty\)，得到

\[
\|f-I_nf\|_\infty
\le(1+\Lambda_n)
\inf_{q\in P_n}
\|f-q\|_\infty.
\]

这条不等式把插值误差拆成两件独立的事。第一件——\(\inf_{q\in P_n}\|f-q\|_\infty\)——是"\(n\) 次多项式空间本身能在多大程度上逼近 \(f\)"，这是多项式空间的能力问题。第二件——\((1+\Lambda_n)\)——是"你选的插值算子会把最佳逼近误差放大多少倍"，这是节点选择的条件数问题。

Runge 现象反映的是第二件：等距高次插值的算子越来越病态，\(\Lambda_n\) 指数增长，把本可以很小的逼近误差放大了成千上万倍。它不否定多项式空间逼近光滑函数的能力——对那个温顺的 \(1/(1+25x^2)\)，\(n\) 次最佳一致逼近多项式的误差确实随 \(n\) 增大而减小。

\subsection{把节点推向端点：Chebyshev 节点}

等距节点的问题出在端点，所以修正也在端点。回看上一节的误差公式：

\[
f(x)-p_n(x)
=\frac{f^{(n+1)}(\xi_x)}{(n+1)!}
\prod_{i=0}^n(x-x_i).
\]

导数因子由 \(f\) 决定，我们控制不了。节点乘积 \(\prod(x-x_i)\) 完全由节点位置决定——这就是我们可以优化的手柄。Chebyshev 节点的核心思想极其简单：选择节点位置，让这个乘积在区间 \([-1,1]\) 上的最大绝对值尽可能小。

在 \([-1,1]\) 上，满足这个极小化条件的节点是 Chebyshev 多项式的根：

\[
x_j=\cos\frac{(2j+1)\pi}{2n+2},
\qquad j=0,\ldots,n.
\]

这些节点在端点附近密集、在区间中心稀疏——恰好与等距节点相反。端点更需要控制，因为高次多项式在区间端点天然更容易放大，等距布点没有为这种几何不均匀性预留任何调控；Chebyshev 节点用密度分布消解了这种不均匀。

用 Chebyshev 节点时，Lebesgue 常数的增长从指数级降为对数级：\(\Lambda_n=O(\log n)\)。对 Runge 函数做 Chebyshev 插值，误差随 \(n\) 单调下降，到 \(n=100\) 时整个区间上几乎不可见任何振荡。对解析函数，Chebyshev 插值常呈现"谱收敛"——误差随 \(n\) 以几何速率下降，快于任何固定幂次的代数收敛。

为什么端点需要更多点，这里有一个更几何化的说法：在一个固定区间上，高次多项式在中心区域可以平缓变化，靠近边界时为了同时满足"在区间内取合适值"和"在区间边界结束"两个约束，弯曲度会急剧上升。更多靠近边界的节点就像给弹簧加了更多固定支点，防止它弹飞。

\subsection{好节点和好算法是两笔独立的账}

第一节介绍的 barycentric 公式能让给定节点上的插值求值更稳定高效，但如果你用的是等距节点，barycentric 公式也无能为力——它解决的是"给定节点后如何稳定求值"的算法问题，不碰"选什么节点"的条件数问题。稳定的算法不能消除一个病态的问题。这与前一单元反复强调的条件数与稳定性的区分完全一致。

反过来，Chebyshev 节点把问题条件数从指数级降为对数级，但如果求值实现还在用幂基 Vandermonde 系统硬解——那套矩阵条件数对几乎任何节点都糟——你就把自己花在节点选择上的一番功夫浪费了。好节点保证问题本身不会被微小扰动放大；好算法保证求解过程不会额外引进扰动。两张牌必须同时握在手里。

\subsection{Runge 不是所有振荡的统称}

Runge 现象具体而特殊，它不应该被泛化为"高次插值就是容易振荡"。下面是几个容易混淆但性质不同的振荡来源：

\begin{itemize}
\tightlist
\item
  并非所有函数的等距插值都发散。低次插值、某些特殊函数（如指数函数在适当区间上）的等距插值表现良好。Runge 是一个反例，说明"可以"发散，不是"必然"发散。
\item
  Chebyshev 节点不能消除 Gibbs 现象。Gibbs 现象出现在用连续函数（如 Fourier 级数部分和或多项式插值）逼近有跳跃间断的函数时——无论用什么节点，跳跃附近的超调不会消失，这是不连续本身的代价。Runge 和 Gibbs 的病因完全不同。
\item
  数据含噪声时，高次精确插值会把噪声一并吞进模型。这不再是逼近论问题，而是统计与模型选择问题——你需要的也许不是插值，而是最小二乘或正则化回归。
\item
  区间外外推的危险远大于区间内插值。即使区间内插值误差在 \(10^{-10}\) 量级，外推到区间外 10\% 的位置，误差可能已经超过函数值本身。外推时你失去了"两侧有约束"的结构，多项式走出最后一个节点后就只剩下"惯性的延伸"。
\end{itemize}

如果数据点位置无法改变——比如来自固定采样率的传感器——可以考虑分段低次插值、样条（本单元最后一节会详细展开）、最小二乘（下一节）或有理函数逼近。不要把"增加全局多项式次数"当作唯一的精度旋钮。旋钮不止一个，而且这个旋钮拧过头会炸。

\subsection{检查精度不要只盯着数据点}

假设有人展示一个 30 次等距插值，并告诉你"看，所有给定数据点上的残差都小于 \(10^{-14}\)，所以这个模型极其精确"。这句话毫无说服力——插值多项式本来就被构造来精确穿过这些数据点，数据点上的零残差是输入条件，不是输出成就。

判断一个插值模型是否可靠，至少需要做三件事：

第一，在整个区间内另外取一批检查点（不用于构造插值的点），看误差在哪些区域突然升高。如果端点附近的误差是中心区域的万倍以上，你手上的就是等距高次插值的典型症状。

第二，轻微扰动原始数据——比如在每个 \(y_j\) 上加一个 \(10^{-10}\) 量级的随机扰动——重新插值，看输出的多项式变化有多大。变化剧烈，说明问题条件数很差，原始数据的任何测量误差都会被灾难性放大。

第三，和其他方案比较——Chebyshev 节点插值、分段低次插值、或样条插值。如果更简单的分段方法已经给出视觉上不可分辨的结果，并且没有端点摆动，那高次全局插值方案就没有理由被保留。

只有数据点上的误差很小，不能证明整个区间上的逼近是准确的。一个在所有证据点上都准确的模型可能是一个在证据点之间乱窜的模型——而且越复杂的模型，乱窜的空间越大。
